\chapter{Literature Review}
\label{chap:background}

This work is based on previous researches in the field of computer science,
including ones on language server and compiler construction. In this
chapter we conduct the literature review of the most relevant sources in these areas.

\section{Language Server Protocol}

The Language Server Protocol (LSP) is a language-agnostic protocol designed to
ease the integration of different programming languages with development
environments. Before the introduction of LSP, language developers had to 
implement and maintain distinct integration for each supported IDE, which was
both time-consuming and resource-intensive. LSP offers a unified protocol that
could be implemented once per language and used across multiple editors.

The current specification of LSP, version 3.17
\cite{language-server-protocol:2016}, defines a standard suite of core editing
features, such as `go to definition', `symbol renaming', `code completion', and
`semantic highlighting'. These features improve the development experience by
providing essential tools for code navigation. 

It is worth noting, that, LSP was designed as an extensible protocol, allowing
language developers to implement additional functionalities via custom requests
and responses.

As an example of such custom functionality, the OCaml language
server~\cite{opam-ocaml-lsp} introduces a custom request method,
$\texttt{ocamllsp/switchImplIntf}$, which allowes the developer to toggle
between interface ($\texttt{foo.mli}$) and implementation ($\texttt{foo.ml}$)
files. This feature highlights the flexibility of LSP to accommodate
language-specific requirements.


\subsection{Language server implementations}
\label{sec:background:implementations}

For optimal integration with the Jasmin compiler, which is implemented in OCaml,
the language server was developed in OCaml too. This choice provides direct
access to compiler API for essential operations, including source code parsing
and functions on typed AST. Below, we analyze relevant OCaml implementations to
derive best practices for language server design.  

\subsubsection*{Jasmin language server}

Aside from the implementation discussed in this work, the Jasmin Language Server
described in \cite{gh-jasmin-lsp} is the only existing option for providing
basic diagnostics support for Jasmin. While functional, this implementation has
several limitations:

\begin{itemize}
    \item It uses a custom-written Language Server Protocol (LSP) wrapper.
    This choice complicates maintenance and extensibility, especially given the
    availability of community-supported LSP libraries for OCaml, such as
    \cite{opam-ocaml-lsp}.
    
    \item Logging is performed directly to \texttt{stderr}, which is
    non-extensible and limits the ability to effectively debug and trace issues.
    A more robust logging framework would facilitate better development and
    maintenance processes.
    
    \item The server currently supports only the
    $\texttt{textDocument/publishDiagnostics}$ functionality. It omits other
    essential LSP features such as hover information, go-to definition, and find
    references, which are critical for improving the user experience and
    productivity.
    
    \item The server lacks configurability, meaning it cannot parameterize
    settings such as the include path or message severity levels. This
    restriction may limit its applicability in some projects.

\end{itemize}

\subsubsection*{OCaml language server (Merlin)}

Another example in the field of language servers is Merlin, a sophisticated tool
specifically designed for the OCaml programming language.  Merlin's development
started in 2013, long before the introduction of the Language Server Protocol
(LSP). In their report \cite{merlin:2018}, the developers and researchers behind
Merlin share valuable insights and methodologies that were instrumental in
crafting this language server \cite{gh-merlin}. 

One of the major contributions of the Merlin project is its sophisticated
management of incomplete code during real-time editing, a feature that is
essential to providing precise diagnostics and stable code suggestions, even when
the code is incomplete. This ability is crucial for developers as it supports a
seamless programming experience by detecting errors at the earliest, providing
real-time autocomplete suggestions. Becoming able to do this is is challenging
and needs a powerful parsing technique and strong type inference mechanisms for
developing well-informed assumptions of the programmer's intentions, despite
incomplete input.

Moreover, Merlin's architecture shows the potential to integrate deeply with
IDEs, offering features like type checking, code navigation, and more. These
capabilities not only enhance the development experience but also set a high
standard for language servers in terms of user support. 

\subsubsection*{Semgrep language server}

Particularly interesting language server implementation is the Semgrep Language Server
\cite{gh-semgrep-lsp}, which served as a primary technical reference due to its
clarity and modular design. 

While the implementation supports only the core LSP features, it presents
pragmatic architectural choices that help maintaining efficiency and simplicity
of a language server. For example, it uses the \texttt{ocaml-lsp}
library~\cite{opam-ocaml-lsp} for LSP messages handling, which simplifies
communication between the editor and the server by abstracting away the
complexities of the LSP implementation. Furthermore, the Semgrep Language Server
utilizes the \texttt{lwt} library \cite{opam-lwt} for asynchronous input/output
handling, enabling non-blocking operations that can be run in a separate thread.

Additionally, the modular design of Semgrep language server helps with 
new feature addition, due to the fact that most of Semgrep language server 
internal functions are implemented as pure functions, that greatly helps 
with debugging and testing new features.

\subsection{Libraries for language server development}

A widely adopted OCaml library for language server development is the
aforementioned \texttt{ocaml-lsp} library \cite{opam-ocaml-lsp}. As of this
writing, this library provides comprehensive support for processing JSON-RPC 2.0
\cite{jsonrpc:2010} communication, which is the underlying protocol for the
Language Server Protocol (LSP). 

One of the key features of the \texttt{ocaml-lsp} library is its capability to
handle bidirectional conversion between raw protocol messages and typed OCaml
representations of LSP structures. This allows developers to work with
strongly-typed data, reducing the likelihood of runtime errors.

Internally, the library relies on the \texttt{yojson} library \cite{gh-yojson}
for efficient JSON parsing and serialization. \texttt{yojson} is well-suited for
this task due to its performance and ease of use, enabling rapid conversion
between JSON data and OCaml types. 

The architecture of \texttt{ocaml-lsp} enables developers to concentrate on
implementing language-specific logic, such as syntax analysis, code completion,
and other higher-order functionalities, without delving into the low-level
details of protocol handling. Furthermore, the library's open-source nature
invites community contributions, leading to continuous improvements and feature
enhancements.

\section{Jasmin Compiler}

The Jasmin compiler \cite{jasmin:2017} is a formally verified toolchain
specifically designed for the compilation of high-assurance cryptographic code
into efficient and secure assembly. One of its primary features is the ability
to verify program properties such as memory safety and termination through its
integrated static analyzer. This ensures that the generated assembly code is not
only performant but also compliant to security standards, making it highly
suitable for cryptographic applications. However, utilizing the full range of
the compiler's utilities necessitates a deep understanding of its source code,
as there is no public API provided for more straightforward interaction.

A distinctive attribute of the Jasmin compiler is its capability to extract
programs to EasyCrypt \cite{easycrypt:2013} specifications. This feature is
particularly significant for advancing the formal verification of cryptographic
algorithms, as it allows developers to conduct thorough correctness and security
proofs within the EasyCrypt framework. The authors assert that for a Jasmin
program that meets safety criteria, its corresponding EasyCrypt embedding
remains sound. This claim underscores the integrated approach of Jasmin and
EasyCrypt in ensuring the verification of cryptographic software, highlighting
the importance of formal methods in the development of secure and verified
systems. Such a combination not only enhances the assurance of security
properties but also facilitates ongoing formal analysis and proofs of modern
cryptographic algorithms.

\section{Tree-sitter}

Tree-sitter is a powerful and flexible parser generator tool designed to build
fast and robust parsers that produce concrete syntax trees (CST) for use in code
editors and other tools. Developed by Max Brunsfeld and first released in 2018,
Tree-sitter has quickly gained popularity due to its efficiency, ease of use,
and comprehensive support for a wide variety of programming languages
\cite{tree-sitter}. 

One of the primary advantages of Tree-sitter over traditional parsing techniques
is its incremental parsing capability, which allows parsers to update the syntax
tree efficiently in response to code changes. This feature is particularly
beneficial in the context of text editors and Integrated Development
Environments. Incremental parsing reduces the computational overhead
associated with re-parsing the entire file, enabling smoother and more
intuitive coding experiences for developers.
